\begin{figure}[H]
  \centering
  \begin{tikzpicture}[                                        
      insignia/.style={rounded corners=4pt, text=white, align=center,
          font=\small\bfseries, minimum width=2.7cm,
          minimum height=0.9cm, inner sep=3pt},
      subt/.style={font=\footnotesize\itshape, align=center, text=black},
      bando/.style={rounded corners=6pt, draw=diagGrupoBorde, fill=diagGrupo,
          inner sep=12pt}
    ]
    % --- Bando izquierdo: modelos propietarios ---
    \node[insignia, fill=marcaOpenAI] (gpt)    at (0,3.1) {GPT\\[-1pt]{\footnotesize\mdseries OpenAI}};
    \node[insignia, fill=marcaClaude] (claude) at (0,1.8) {Claude\\[-1pt]{\footnotesize\mdseries Anthropic}};
    \node[insignia, fill=marcaGemini] (gemini) at (0,0.5) {Gemini\\[-1pt]{\footnotesize\mdseries Google}};
    \node[subt] (subP) at (0,-0.9) {Pesos cerrados\\Acceso mediante API};
    % --- Bando derecho: modelos de pesos abiertos ---
    \node[insignia, fill=marcaLlama]   (llama)   at (7,3.1) {Llama\\[-1pt]{\footnotesize\mdseries Meta}};
    \node[insignia, fill=marcaMistral] (mistral) at (7,1.8) {Mistral\\[-1pt]{\footnotesize\mdseries Mistral AI}};
    \node[insignia, fill=marcaQwen]    (qwen)    at (7,0.5) {Qwen\\[-1pt]{\footnotesize\mdseries Alibaba}};
    \node[subt] (subA) at (7,-0.9) {Parámetros descargables\\Ejecución local · control de los datos};
    % --- VS central ---
    \node[circle, draw=diagNodoBorde, fill=white, line width=1pt,
      minimum size=1.25cm, font=\LARGE\bfseries, text=diagNodoBorde] (vs) at (3.0,1.8) {VS};
    % Paneles de fondo con su título
    \begin{scope}[on background layer]
      \node[bando, fit=(gpt)(gemini)(subP),
        label={[font=\bfseries]north:Modelos propietarios}] {};
      \node[bando, fit=(llama)(qwen)(subA),
        label={[font=\bfseries]north:Modelos de pesos abiertos}] {};
    \end{scope}
  \end{tikzpicture}
  \caption{Modelos de lenguaje propietarios (pesos
    cerrados, acceso mediante la API del proveedor) frente a modelos de pesos
    abiertos (parámetros descargables, ejecución local y control de los datos).
    Los nombres se muestran como insignias identificativas, no como los
    logotipos oficiales de cada marca. Fuente: elaboración propia.}
  \label{fig:privativos_vs_abiertos}
\end{figure}
